\documentclass[11pt]{article}
\usepackage{amsmath, amssymb, amsthm}
\usepackage{geometry}
\geometry{letterpaper, margin=1in}
\usepackage{hyperref}

\title{Math - Test Sections Questions: \\ Nonhomogeneous Equations: The Method of Undetermined Coefficients}
\author{Curated Problem Sets}
\date{}

\begin{document}

\maketitle
\tableofcontents
\newpage

% ---------------------------------------------------------
\section{Concept Checklist \& Question Mapping}
% ---------------------------------------------------------
Before diving into the problems, ensure you understand the following concepts and problem types derived from the text. The problems in this document are mapped to these specific concepts to ensure comprehensive coverage.

\begin{itemize}
    \item \textbf{C1: Applicability Checks} - Determining if the method can be applied based on the form of $f(t)$. (Problems 1, 2, 3)
    \item \textbf{C2: Basic Undetermined Coefficients} - Standard workflow where the guess does not overlap with the homogeneous solution ($s=0$). (Problems 4, 5, 6, 7)
    \item \textbf{C3: Root Overlap / Resonance (Simple Root)} - The guess overlaps with a simple root of the characteristic equation, requiring a $t^1$ multiplier ($s=1$). (Problems 8, 9, 10, 11)
    \item \textbf{C4: Root Overlap / Resonance (Double Root)} - The guess overlaps with a double root of the characteristic equation, requiring a $t^2$ multiplier ($s=2$). (Problems 12, 13, 14)
    \item \textbf{C5: Form Setup Only} - Constructing complex trial functions without evaluating the coefficients. (Problems 15, 16, 17, 18)
    \item \textbf{C6: Superposition of Multiple Terms} - Handling $f(t)$ that is a sum of different function types by splitting the guess. (Problems 19, 20, 21)
    \item \textbf{C7: Products of Three Function Types} - Handling $f(t)$ containing a polynomial, exponential, and trigonometric function simultaneously. (Problems 22, 23, 24)
    \item \textbf{C8: Higher-Order Equations} - Applying the method to 3rd-order or higher differential equations. (Problems 25, 26, 27)
    \item \textbf{C9: Diagnostic Testing / "Find the Flaw"} - Identifying common traps such as the "Missing Sibling", "Lazy Polynomial", or missed multipliers. (Problems 28, 29, 30)
\end{itemize}

\newpage

% ---------------------------------------------------------
\section{Problem Sets}
% ---------------------------------------------------------

\subsection*{Applicability Checks}
\textbf{1.} Decide whether or not the method of undetermined coefficients can be applied to $y'' - 4y' + 4y = t^{-2}e^{2t}$. \\
\textbf{2.} Decide whether or not the method of undetermined coefficients can be applied to $y'' + 9y = \tan(3t)$. \\
\textbf{3.} Decide whether or not the method of undetermined coefficients can be applied to $y'' - y = t^4 \sin(2t) e^{-t}$.

\subsection*{Basic Undetermined Coefficients}
\textbf{4.} Find a particular solution to $y'' + 4y = 12$. \\
\textbf{5.} Find a particular solution to $y'' - 3y' + 2y = 4e^{3t}$. \\
\textbf{6.} Find a particular solution to $y'' + 2y' + 5y = 10\sin t$. \\
\textbf{7.} Find a particular solution to $y'' - y' - 2y = 2x^2$.

\subsection*{Root Overlap / Resonance (Simple Root)}
\textbf{8.} Find a particular solution to $y'' - 5y' + 4y = 3e^{4t}$. \\
\textbf{9.} Find a particular solution to $y'' + 9y = 6\cos(3t)$. \\
\textbf{10.} Find a particular solution to $y'' + 2y' = 4t$. \\
\textbf{11.} Find a particular solution to $y'' - y = 2te^t$.

\subsection*{Root Overlap / Resonance (Double Root)}
\textbf{12.} Find a particular solution to $y'' - 6y' + 9y = 2e^{3t}$. \\
\textbf{13.} Find a particular solution to $y'' + 4y' + 4y = 12te^{-2t}$. \\
\textbf{14.} Find a particular solution to $y'' = 6t$.

\subsection*{Form Setup Only (Do not evaluate coefficients)}
\textbf{15.} Determine the form of a particular solution for $y'' + 25y = t^2 \cos(5t)$. \\
\textbf{16.} Determine the form of a particular solution for $y'' - 8y' + 16y = t^3 e^{4t}$. \\
\textbf{17.} Determine the form of a particular solution for $y'' - 2y' + 5y = e^t \sin(2t)$. \\
\textbf{18.} Determine the form of a particular solution for $y'' - 7y' + 10y = t^2 e^{3t} \cos(t)$.

\subsection*{Superposition of Multiple Terms}
\textbf{19.} Find a particular solution to $y'' - y = 2t + 4e^{3t}$. \\
\textbf{20.} Find a particular solution to $y'' + 4y = 8 + \sin t$. \\
\textbf{21.} Determine the form of a particular solution for $y'' - 3y' + 2y = e^t + t\sin t$.

\subsection*{Products of Three Function Types}
\textbf{22.} Determine the form of a particular solution for $y'' + y = t e^{2t} \cos t$. \\
\textbf{23.} Determine the form of a particular solution for $y'' - 4y' + 5y = t^2 e^{2t} \sin t$. \\
\textbf{24.} Find a particular solution to $y'' - 2y' + y = e^t \cos t$.

\subsection*{Higher-Order Equations}
\textbf{25.} Find a particular solution to $y''' - y'' = 2$. \\
\textbf{26.} Determine the form of a particular solution for $y^{(4)} - y = e^t$. \\
\textbf{27.} Find a particular solution to $y''' + y' = \sin t$.

\subsection*{Diagnostic Testing: "Find the Flaw"}
Identify the conceptual or algebraic flaw in the following proposed solutions, explain why it is wrong in one sentence, and provide the correct setup. \\

\noindent \textbf{28. Flawed Solution A} \\
\textit{Problem:} Find a particular solution to $y'' + 4y = \sin(2t)$. \\
\textit{Solution:} Since the right side is $\sin(2t)$, we guess $y_p(t) = A\cos(2t) + B\sin(2t)$. \\

\noindent \textbf{29. Flawed Solution B} \\
\textit{Problem:} Find a particular solution to $y'' - 2y' + y = t e^t$. \\
\textit{Solution:} The characteristic equation has a double root at $r=1$. The nonhomogeneous term is $te^t$. Since $r=1$ is a double root, we must multiply our guess by $t^2$. We guess $y_p(t) = (Ate^t) \cdot t^2 = At^3e^t$. \\

\noindent \textbf{30. Flawed Solution C} \\
\textit{Problem:} Find a particular solution to $y'' + y = \ln t$. \\
\textit{Solution:} The nonhomogeneous term is $\ln t$. We guess a logarithmic function of the same form: $y_p(t) = A\ln t + B$.

\newpage

% ---------------------------------------------------------
\section{Detailed Solutions}
% ---------------------------------------------------------

\subsection*{Applicability Checks}

\textbf{Solution 1:} 
The nonhomogeneous term is $f(t) = t^{-2}e^{2t}$. The method of undetermined coefficients strictly requires the polynomial component to have non-negative integer powers. Because $f(t)$ contains $t^{-2}$, the method \textbf{cannot be applied}.

\textbf{Solution 2:} 
The nonhomogeneous term is $f(t) = \tan(3t)$. The method only applies to polynomials, exponentials, sines, cosines, or finite products/sums of these. It does not apply to tangent functions. Thus, the method \textbf{cannot be applied}.

\textbf{Solution 3:} 
The nonhomogeneous term is $f(t) = t^4 \sin(2t) e^{-t}$. This is a product of a polynomial with a non-negative integer power ($t^4$), a sine function, and an exponential function. This perfectly matches the standard form. Thus, the method \textbf{can be applied}.

\subsection*{Basic Undetermined Coefficients}

\textbf{Solution 4:} 
Characteristic equation: $r^2 + 4 = 0 \implies r = \pm 2i$. 
$f(t) = 12$. Here $\alpha=0, \beta=0$. Since $0$ is not a root, $s=0$. 
Guess: $y_p(t) = A$. 
Differentiating: $y_p' = 0, y_p'' = 0$. 
Substitute: $0 + 4A = 12 \implies A = 3$. 
\textbf{Answer:} $y_p(t) = 3$.

\textbf{Solution 5:} 
Characteristic equation: $r^2 - 3r + 2 = 0 \implies (r-1)(r-2) = 0 \implies r = 1, 2$. 
$f(t) = 4e^{3t}$. Here $\alpha=3$. Since $3$ is not a root, $s=0$. 
Guess: $y_p(t) = Ae^{3t}$. 
Differentiating: $y_p' = 3Ae^{3t}, y_p'' = 9Ae^{3t}$. 
Substitute: $9Ae^{3t} - 3(3Ae^{3t}) + 2(Ae^{3t}) = 4e^{3t} \implies 2Ae^{3t} = 4e^{3t} \implies 2A = 4 \implies A = 2$. 
\textbf{Answer:} $y_p(t) = 2e^{3t}$.

\textbf{Solution 6:} 
Characteristic equation: $r^2 + 2r + 5 = 0 \implies r = -1 \pm 2i$. 
$f(t) = 10\sin t$. Here $\alpha=0, \beta=1 \implies \alpha+i\beta = i$. Since $i$ is not a root, $s=0$. 
Guess: $y_p(t) = A\cos t + B\sin t$. 
$y_p' = -A\sin t + B\cos t$ \\
$y_p'' = -A\cos t - B\sin t$ \\
Substitute: 
$(-A\cos t - B\sin t) + 2(-A\sin t + B\cos t) + 5(A\cos t + B\sin t) = 10\sin t$. \\
Group cosine terms: $-A + 2B + 5A = 4A + 2B = 0 \implies B = -2A$. \\
Group sine terms: $-B - 2A + 5B = -2A + 4B = 10$. \\
Substitute $B$: $-2A + 4(-2A) = 10 \implies -10A = 10 \implies A = -1$. Then $B = 2$. \\
\textbf{Answer:} $y_p(t) = -\cos t + 2\sin t$.

\textbf{Solution 7:} 
Characteristic equation: $r^2 - r - 2 = 0 \implies r = 2, -1$. 
$f(x) = 2x^2$. Here $\alpha=0$. Since $0$ is not a root, $s=0$. 
Guess: $y_p(x) = Ax^2 + Bx + C$. 
$y_p' = 2Ax + B$, $y_p'' = 2A$. 
Substitute: $2A - (2Ax + B) - 2(Ax^2 + Bx + C) = 2x^2$. \\
$-2Ax^2 + (-2A - 2B)x + (2A - B - 2C) = 2x^2$. \\
Equate coefficients: \\
$x^2$: $-2A = 2 \implies A = -1$. \\
$x$: $-2A - 2B = 0 \implies 2 - 2B = 0 \implies B = 1$. \\
Constant: $2A - B - 2C = 0 \implies -2 - 1 - 2C = 0 \implies 2C = -3 \implies C = -3/2$. \\
\textbf{Answer:} $y_p(x) = -x^2 + x - \frac{3}{2}$.

\subsection*{Root Overlap / Resonance (Simple Root)}

\textbf{Solution 8:} 
Characteristic equation: $r^2 - 5r + 4 = 0 \implies r = 1, 4$. 
$f(t) = 3e^{4t}$. Here $\alpha=4$. Since $4$ is a simple root, $s=1$. 
Guess: $y_p(t) = Ate^{4t}$. 
$y_p' = A(4t+1)e^{4t}$, $y_p'' = A(16t+8)e^{4t}$. 
Substitute: $A(16t+8)e^{4t} - 5A(4t+1)e^{4t} + 4Ate^{4t} = 3e^{4t}$. \\
Divide by $e^{4t}$: $16At + 8A - 20At - 5A + 4At = 3$. \\
Notice the $t$ terms cancel: $(16A - 20A + 4A)t + 3A = 3 \implies 3A = 3 \implies A = 1$. \\
\textbf{Answer:} $y_p(t) = te^{4t}$.

\textbf{Solution 9:} 
Characteristic equation: $r^2 + 9 = 0 \implies r = \pm 3i$. 
$f(t) = 6\cos(3t)$. Here $\alpha=0, \beta=3 \implies 3i$. Since $3i$ is a simple root, $s=1$. 
Guess: $y_p(t) = t(A\cos 3t + B\sin 3t)$. 
$y_p' = (A\cos 3t + B\sin 3t) + t(-3A\sin 3t + 3B\cos 3t)$. \\
$y_p'' = (-3A\sin 3t + 3B\cos 3t) + (-3A\sin 3t + 3B\cos 3t) + t(-9A\cos 3t - 9B\sin 3t)$ \\
$y_p'' = -6A\sin 3t + 6B\cos 3t - 9t(A\cos 3t + B\sin 3t)$. \\
Substitute: $[-6A\sin 3t + 6B\cos 3t - 9y_p] + 9y_p = 6\cos 3t$. \\
$-6A\sin 3t + 6B\cos 3t = 6\cos 3t$. \\
Equate coefficients: $-6A = 0 \implies A = 0$. $6B = 6 \implies B = 1$. \\
\textbf{Answer:} $y_p(t) = t\sin 3t$.

\textbf{Solution 10:} 
Characteristic equation: $r^2 + 2r = 0 \implies r = 0, -2$. 
$f(t) = 4t$. Here $\alpha=0$. Since $0$ is a simple root, $s=1$. 
Guess: $y_p(t) = t(At + B) = At^2 + Bt$. 
$y_p' = 2At + B$, $y_p'' = 2A$. 
Substitute: $2A + 2(2At + B) = 4t \implies 4At + (2A + 2B) = 4t$. \\
Equate coefficients: $4A = 4 \implies A = 1$. \\
$2A + 2B = 0 \implies 2(1) + 2B = 0 \implies B = -1$. \\
\textbf{Answer:} $y_p(t) = t^2 - t$.

\textbf{Solution 11:} 
Characteristic equation: $r^2 - 1 = 0 \implies r = \pm 1$. 
$f(t) = 2te^t$. Here $\alpha=1$. Since $1$ is a simple root, $s=1$. 
Guess: $y_p(t) = t(At + B)e^t = (At^2 + Bt)e^t$. 
$y_p' = (2At + B)e^t + (At^2 + Bt)e^t = (At^2 + (2A+B)t + B)e^t$. \\
$y_p'' = (2At + 2A + B)e^t + (At^2 + (2A+B)t + B)e^t = (At^2 + (4A+B)t + 2A + 2B)e^t$. \\
Substitute: $(At^2 + (4A+B)t + 2A + 2B)e^t - (At^2 + Bt)e^t = 2te^t$. \\
Divide by $e^t$: $4At + 2A + 2B = 2t$. \\
Equate coefficients: $4A = 2 \implies A = 1/2$. \\
$2A + 2B = 0 \implies 1 + 2B = 0 \implies B = -1/2$. \\
\textbf{Answer:} $y_p(t) = (\frac{1}{2}t^2 - \frac{1}{2}t)e^t$.

\subsection*{Root Overlap / Resonance (Double Root)}

\textbf{Solution 12:} 
Characteristic equation: $r^2 - 6r + 9 = 0 \implies (r-3)^2 = 0 \implies r = 3$ (double root). 
$f(t) = 2e^{3t}$. Here $\alpha=3$. Since $3$ is a double root, $s=2$. 
Guess: $y_p(t) = At^2e^{3t}$. 
$y_p' = A(3t^2 + 2t)e^{3t}$. \\
$y_p'' = A(9t^2 + 6t + 6t + 2)e^{3t} = A(9t^2 + 12t + 2)e^{3t}$. \\
Substitute: $A(9t^2 + 12t + 2) - 6A(3t^2 + 2t) + 9At^2 = 2$. \\
$t^2(9A - 18A + 9A) + t(12A - 12A) + 2A = 2 \implies 2A = 2 \implies A = 1$. \\
\textbf{Answer:} $y_p(t) = t^2e^{3t}$.

\textbf{Solution 13:} 
Characteristic equation: $r^2 + 4r + 4 = 0 \implies (r+2)^2 = 0 \implies r = -2$ (double root). 
$f(t) = 12te^{-2t}$. Here $\alpha=-2$. Since $-2$ is a double root, $s=2$. 
Guess: $y_p(t) = t^2(At + B)e^{-2t} = (At^3 + Bt^2)e^{-2t}$. 
$y_p' = (3At^2 + 2Bt)e^{-2t} - 2(At^3 + Bt^2)e^{-2t} = (-2At^3 + (3A-2B)t^2 + 2Bt)e^{-2t}$. \\
$y_p'' = (-6At^2 + 2(3A-2B)t + 2B)e^{-2t} - 2(-2At^3 + (3A-2B)t^2 + 2Bt)e^{-2t}$ \\
$y_p'' = (4At^3 + (-12A+4B)t^2 + (6A-8B)t + 2B)e^{-2t}$. \\
Substitute into $y'' + 4y' + 4y$: \\
$t^3[4A - 8A + 4A] + t^2[-12A+4B + 12A-8B + 4B] + t[6A-8B + 8B] + 2B = 12t$. \\
$6At + 2B = 12t$. \\
Equate coefficients: $6A = 12 \implies A = 2$. $2B = 0 \implies B = 0$. \\
\textbf{Answer:} $y_p(t) = 2t^3e^{-2t}$.

\textbf{Solution 14:} 
Characteristic equation: $r^2 = 0 \implies r = 0$ (double root). 
$f(t) = 6t$. Here $\alpha=0$. Since $0$ is a double root, $s=2$. 
Guess: $y_p(t) = t^2(At + B) = At^3 + Bt^2$. 
$y_p' = 3At^2 + 2Bt$, $y_p'' = 6At + 2B$. 
Substitute: $6At + 2B = 6t$. \\
Equate coefficients: $6A = 6 \implies A = 1$. $2B = 0 \implies B = 0$. \\
\textbf{Answer:} $y_p(t) = t^3$.

\subsection*{Form Setup Only}

\textbf{Solution 15:} 
Characteristic equation: $r^2 + 25 = 0 \implies r = \pm 5i$. 
$f(t) = t^2 \cos(5t) \implies \alpha=0, \beta=5 \implies 5i$. Since $5i$ is a simple root, $s=1$. The polynomial is degree 2. \\
\textbf{Answer:} $y_p(t) = t \left( (A_2t^2 + A_1t + A_0)\cos(5t) + (B_2t^2 + B_1t + B_0)\sin(5t) \right)$.

\textbf{Solution 16:} 
Characteristic equation: $r^2 - 8r + 16 = 0 \implies (r-4)^2 = 0 \implies r = 4$ (double root). 
$f(t) = t^3 e^{4t} \implies \alpha=4$. Since $4$ is a double root, $s=2$. The polynomial is degree 3. \\
\textbf{Answer:} $y_p(t) = t^2 (A_3t^3 + A_2t^2 + A_1t + A_0)e^{4t}$.

\textbf{Solution 17:} 
Characteristic equation: $r^2 - 2r + 5 = 0 \implies r = 1 \pm 2i$. 
$f(t) = e^t \sin(2t) \implies \alpha=1, \beta=2 \implies 1+2i$. Since $1+2i$ is a simple root, $s=1$. The polynomial is degree 0. \\
\textbf{Answer:} $y_p(t) = t e^t (A\cos(2t) + B\sin(2t))$.

\textbf{Solution 18:} 
Characteristic equation: $r^2 - 7r + 10 = 0 \implies r = 2, 5$. 
$f(t) = t^2 e^{3t} \cos(t) \implies \alpha=3, \beta=1 \implies 3+i$. Since $3+i$ is not a root, $s=0$. The polynomial is degree 2. \\
\textbf{Answer:} $y_p(t) = e^{3t} \left( (A_2t^2 + A_1t + A_0)\cos t + (B_2t^2 + B_1t + B_0)\sin t \right)$.

\subsection*{Superposition of Multiple Terms}

\textbf{Solution 19:} 
Characteristic equation: $r^2 - 1 = 0 \implies r = \pm 1$. 
Split $f(t) = 2t + 4e^{3t}$ into two problems. \\
Part 1: $y_{p1}$ for $2t$. $\alpha=0, s=0$. Guess $y_{p1} = At + B$. $y_{p1}'' = 0$. \\
$0 - (At + B) = 2t \implies A = -2, B = 0 \implies y_{p1} = -2t$. \\
Part 2: $y_{p2}$ for $4e^{3t}$. $\alpha=3, s=0$. Guess $y_{p2} = Ce^{3t}$. $y_{p2}'' = 9Ce^{3t}$. \\
$9Ce^{3t} - Ce^{3t} = 4e^{3t} \implies 8C = 4 \implies C = 1/2 \implies y_{p2} = \frac{1}{2}e^{3t}$. \\
\textbf{Answer:} $y_p(t) = -2t + \frac{1}{2}e^{3t}$.

\textbf{Solution 20:} 
Characteristic equation: $r^2 + 4 = 0 \implies r = \pm 2i$. 
Split $f(t) = 8 + \sin t$. \\
Part 1: $y_{p1}$ for $8$. $\alpha=0, s=0$. Guess $y_{p1} = A$. $4A = 8 \implies A = 2$. \\
Part 2: $y_{p2}$ for $\sin t$. $\alpha=0, \beta=1, s=0$. Guess $y_{p2} = B\cos t + C\sin t$. \\
$y_{p2}'' = -B\cos t - C\sin t$. \\
$(-B\cos t - C\sin t) + 4(B\cos t + C\sin t) = \sin t \implies 3B\cos t + 3C\sin t = \sin t$. \\
$3B = 0 \implies B = 0$. $3C = 1 \implies C = 1/3$. \\
\textbf{Answer:} $y_p(t) = 2 + \frac{1}{3}\sin t$.

\textbf{Solution 21:} 
Characteristic equation: $r^2 - 3r + 2 = 0 \implies r = 1, 2$. 
Split $f(t) = e^t + t\sin t$. \\
Part 1: For $e^t$, $\alpha=1$. Since $1$ is a simple root, $s=1$. Form: $A t e^t$. \\
Part 2: For $t\sin t$, $\alpha=0, \beta=1 \implies i$. Since $i$ is not a root, $s=0$. Form: $(B_1t + B_0)\cos t + (C_1t + C_0)\sin t$. \\
\textbf{Answer:} $y_p(t) = A t e^t + (B_1t + B_0)\cos t + (C_1t + C_0)\sin t$.

\subsection*{Products of Three Function Types}

\textbf{Solution 22:} 
Characteristic equation: $r^2 + 1 = 0 \implies r = \pm i$. 
$f(t) = t e^{2t} \cos t \implies \alpha=2, \beta=1 \implies 2+i$. Since $2+i$ is not a root, $s=0$. The polynomial is degree 1. \\
\textbf{Answer:} $y_p(t) = e^{2t} \left( (A_1t + A_0)\cos t + (B_1t + B_0)\sin t \right)$.

\textbf{Solution 23:} 
Characteristic equation: $r^2 - 4r + 5 = 0 \implies r = 2 \pm i$. 
$f(t) = t^2 e^{2t} \sin t \implies \alpha=2, \beta=1 \implies 2+i$. Since $2+i$ is a simple root, $s=1$. The polynomial is degree 2. \\
\textbf{Answer:} $y_p(t) = t e^{2t} \left( (A_2t^2 + A_1t + A_0)\cos t + (B_2t^2 + B_1t + B_0)\sin t \right)$.

\textbf{Solution 24:} 
Characteristic equation: $r^2 - 2r + 1 = 0 \implies (r-1)^2 = 0 \implies r = 1$ (double root). 
$f(t) = e^t \cos t \implies \alpha=1, \beta=1 \implies 1+i$. Since $1+i$ is not a root, $s=0$. 
Guess: $y_p(t) = e^t(A\cos t + B\sin t)$. 
$y_p' = e^t(A\cos t + B\sin t - A\sin t + B\cos t) = e^t((A+B)\cos t + (B-A)\sin t)$. \\
$y_p'' = e^t((A+B)\cos t + (B-A)\sin t - (A+B)\sin t + (B-A)\cos t) = e^t(2B\cos t - 2A\sin t)$. \\
Substitute into $y'' - 2y' + y$: \\
$e^t [2B - 2(A+B) + A]\cos t + e^t [-2A - 2(B-A) + B]\sin t = e^t\cos t$. \\
$e^t [-A\cos t - B\sin t] = e^t\cos t$. \\
Equate coefficients: $-A = 1 \implies A = -1$. $-B = 0 \implies B = 0$. \\
\textbf{Answer:} $y_p(t) = -e^t \cos t$.

\subsection*{Higher-Order Equations}

\textbf{Solution 25:} 
Characteristic equation: $r^3 - r^2 = 0 \implies r^2(r-1) = 0 \implies r = 0$ (double), $r = 1$. 
$f(t) = 2 \implies \alpha=0$. Since $0$ is a double root, $s=2$. 
Guess: $y_p(t) = At^2$. 
$y_p' = 2At, y_p'' = 2A, y_p''' = 0$. 
Substitute: $0 - 2A = 2 \implies A = -1$. \\
\textbf{Answer:} $y_p(t) = -t^2$.

\textbf{Solution 26:} 
Characteristic equation: $r^4 - 1 = 0 \implies (r^2-1)(r^2+1) = 0 \implies r = \pm 1, \pm i$. 
$f(t) = e^t \implies \alpha=1$. Since $1$ is a simple root, $s=1$. \\
\textbf{Answer:} $y_p(t) = A t e^t$.

\textbf{Solution 27:} 
Characteristic equation: $r^3 + r = 0 \implies r(r^2+1) = 0 \implies r = 0, \pm i$. 
$f(t) = \sin t \implies \alpha=0, \beta=1 \implies i$. Since $i$ is a simple root, $s=1$. 
Guess: $y_p(t) = t(A\cos t + B\sin t)$. 
$y_p' = (A\cos t + B\sin t) + t(-A\sin t + B\cos t)$. \\
$y_p'' = -2A\sin t + 2B\cos t - t(A\cos t + B\sin t)$. \\
$y_p''' = -3A\cos t - 3B\sin t - t(-A\sin t + B\cos t)$. \\
Substitute into $y''' + y'$: \\
$[-3A\cos t - 3B\sin t - t(-A\sin t + B\cos t)] + [A\cos t + B\sin t + t(-A\sin t + B\cos t)] = \sin t$. \\
$-2A\cos t - 2B\sin t = \sin t$. \\
Equate coefficients: $-2A = 0 \implies A = 0$. $-2B = 1 \implies B = -1/2$. \\
\textbf{Answer:} $y_p(t) = -\frac{1}{2}t\sin t$.

\subsection*{Diagnostic Testing: "Find the Flaw"}

\textbf{Solution 28:} 
\textit{Flaw:} The characteristic roots are $\pm 2i$, which matches the parameters of the nonhomogeneous term ($\alpha=0, \beta=2$), so the standard guess will vanish upon substitution. \\
\textit{Correct Step:} Multiply the guess by $t$: $y_p(t) = t(A\cos(2t) + B\sin(2t))$.

\textbf{Solution 29:} 
\textit{Flaw:} The "Lazy Polynomial" trap: the nonhomogeneous term has a polynomial of degree 1, so the base guess must include a full degree 1 polynomial $(At+B)$ before multiplying by the $t^2$ multiplier. \\
\textit{Correct Step:} Guess $y_p(t) = t^2(At + B)e^t = (At^3 + Bt^2)e^t$.

\textbf{Solution 30:} 
\textit{Flaw:} The method of undetermined coefficients does not apply to logarithmic functions; the failure is conceptual. \\
\textit{Correct Step:} State that the method cannot be used (Variation of Parameters is required).

\newpage

% ---------------------------------------------------------
\section{Concept Checklist \& Question Mapping (Appendix)}
% ---------------------------------------------------------
Use this index to quickly navigate back to specific problem types for review.

\begin{itemize}
    \item \textbf{C1: Applicability Checks} $\rightarrow$ Questions 1, 2, 3
    \item \textbf{C2: Basic Undetermined Coefficients} $\rightarrow$ Questions 4, 5, 6, 7
    \item \textbf{C3: Root Overlap / Resonance (Simple Root)} $\rightarrow$ Questions 8, 9, 10, 11
    \item \textbf{C4: Root Overlap / Resonance (Double Root)} $\rightarrow$ Questions 12, 13, 14
    \item \textbf{C5: Form Setup Only} $\rightarrow$ Questions 15, 16, 17, 18
    \item \textbf{C6: Superposition of Multiple Terms} $\rightarrow$ Questions 19, 20, 21
    \item \textbf{C7: Products of Three Function Types} $\rightarrow$ Questions 22, 23, 24
    \item \textbf{C8: Higher-Order Equations} $\rightarrow$ Questions 25, 26, 27
    \item \textbf{C9: Diagnostic Testing / "Find the Flaw"} $\rightarrow$ Questions 28, 29, 30
\end{itemize}

\end{document}